
\section{Helmholtz Equation for Atmospheric Propagation}

A monochromatic EM source $E(\dim{x},\dim{y},\dim{z}) = A(\dim{x},\dim{y},\dim{z})\exp{-i\omega \dim{t}} + c.c.$, 
propagating in a nondispersive static atmosphere follows a stochastic Helmholtz equation  [Anderson Phillips, "Laser Beam Propagation through random media"]

\begin{equation}\label{EQ:helmholtz_equation_1}
\begin{split}
  &\nabla^2 \dim{A} + k^2  (n_0 +  n(\dim{x},\dim{y},\dim{z}))^2 \dim{A} = 0 \quad  \mbox{for}   
  \quad -\infty < \dim{x},\dim{y} < \infty, \quad 0 < \dim{z} < L, \\
  &\quad  \mbox{where} \quad \dim{A}(\dim{x},\dim{y},0) = \dim{A}_{i}(\dim{x},\dim{y})\exp{i \Phi(\dim{x},\dim{y})},
\end{split}
\end{equation}
and $\dim{A}_{i}(\dim{x},\dim{y})$ is of finite support, \ie $\dim{A}_{i}(\dim{x},\dim{y}) = 0$ 
when $\sqrt{\dim{x}^2 + \dim{y}^2} > D_{i}$. Note, $k = \frac{\omega}{c} = \frac{2\pi}{\lambda}$.  


\subsection{Atmospheric Statistics}
This section contains the most important part of these notes.  Normally, the stochastic part 
of the atmosphere is described by the \textbf{structure function},
\begin{equation}\label{EQ:structure_function_0}
\begin{split}
  \Expect{(n(\vect{\dim{r}}_1) - n(\vect{\dim{r}}_2))^2} &= 
  C_n^2  \begin{cases}
        l_{0}^{-4/3} \|\vect{\dim{r}}_1-\vect{\dim{r}}_2\|^{2} & 0 < \|\vect{\dim{r}}_1-\vect{\dim{r}}_2\| \leq l_{0} \\
       \|\vect{\dim{r}}_1 - \vect{\dim{r}}_2 \|^{2/3} &  l_{0}  \leq \|\vect{\dim{r}}_1 - \vect{\dim{r}}_2\| \leq  L_{0}  \\
   \end{cases}\\
   &= \Expect{n^2(\vect{\dim{r}}_1)} +  \Expect{n^2(\vect{\dim{r}}_2)} - 2\Expect{n(\vect{\dim{r}}_1)n(\vect{\dim{r}}_2)} \\
\end{split}
\end{equation}
Note, this is the only statistic provided to us and this is the only theoretical information we have to act on.
It is common to make the assumption that the index of refraction variations $n(\vect{\dim{r}})$ follow a 
\textbf{stationary} or \textbf{homogeneous} process \href{https://en.wikipedia.org/wiki/Stationary_process}{[Wiki]}, 
at least in the wide-sense~\cite{Wiley_Goodman_1985}, which imposes 
that the $n(\vect{\dim{r}})$ has constant mean and an auto-covariance that only depends on the difference in locations:
 \begin{equation}\label{EQ:structure_function_0}
\begin{split}
  \Expect{n(\vect{\dim{r}}_1)} = \mu_{n}, \quad \Expect{n(\vect{\dim{r}}_1)n(\vect{\dim{r}}_2)} = C_{cos}(\vect{\dim{r}}_1-\vect{\dim{r}}_2):\mathbb{R}^3 \to \mathbb{R}
\end{split}
\end{equation}
In addition, it is commonly assumed that $n(\vect{\dim{r}})$ is mean zero and \textbf{isotropic}, 
implying that $F(\vect{\dim{r}}_1-\vect{\dim{r}}_2) = F(\|\vect{\dim{r}}_1-\vect{\dim{r}}_2\|)$.
This implies that $\Expect{n^2(\vect{\dim{r}}_1)} = \Expect{n^2(\vect{\dim{r}}_2)} = C_{cov}(\| \vect{0} \|) = \sigma_{n}^{2}$ 
and gives  
\begin{equation}\label{EQ:structure_function}
\begin{split}
  \Expect{(n(\vect{\dim{r}}_1) - n(\vect{\dim{r}}_2))^2} &= 
  C_n^2  \begin{cases}
        l_{0}^{-4/3} \|\vect{\dim{r}}_1-\vect{\dim{r}}_2\|^{2} & 0 < \|\vect{\dim{r}}_1-\vect{\dim{r}}_2\| \leq l_{0} \\
       \|\vect{\dim{r}}_1 - \vect{\dim{r}}_2 \|^{2/3} &  l_{0}  \leq \|\vect{\dim{r}}_1 - \vect{\dim{r}}_2\| \leq  L_{0}  \\
   \end{cases}\\
   &= \Expect{n^2(\vect{\dim{r}}_1)} +  \Expect{n^2(\vect{\dim{r}}_2)} - 2\Expect{n(\vect{\dim{r}}_1)n(\vect{\dim{r}}_2)} \\
   &= 2\sigma_{n}^{2}( 1-C_{cor}(\| \vect{\dim{r}}_1 - \vect{\dim{r}}_2\|) )
\end{split}
\end{equation}
where we let $C_{cor}(\cdot) = C_{cov}(\cdot)/\sigma_{n}^{2}$.

so
\begin{equation}\label{EQ:correlation}
  C_{cor}(\|\vect{\dim{r}}_1-\vect{\dim{r}}_{2}\|) = 1 - \frac{C_n^2}{2\sigma_n^2}
    \begin{cases}
        l_{0}^{-4/3} \|\vect{\dim{r}}_1-\vect{\dim{r}}_{2}\|^{2} & 0 < \|\vect{\dim{r}}_1-\vect{\dim{r}}_{2}\| \leq l_{0} \\
        \|\vect{\dim{r}}_1-\vect{\dim{r}}_{2}\|^{2/3} &  l_{0}  \leq \|\vect{\dim{r}}_1-\vect{\dim{r}}_{2}\| \leq  L_{0}  \\
   \end{cases}
\end{equation}

If we assume that $|C_{cor}(\|\vect{\dim{r}}_1-\vect{\dim{r}}_{2}\|)| \ll 1$ for $ \|\vect{\dim{r}}_1-\vect{\dim{r}}_{2}\| \gg L_0$, 
then letting $\|\vect{\dim{r}}_1-\vect{\dim{r}}_{2}\| = C_0 L_0$ for $C_0 > 1$ ($C_0$ is just an adjustment factor) we can write
\begin{equation}
\begin{split}
  0 \approx 1 - C_0^{2/3} \frac{C_n^2 L_0^{2/3}}{2\sigma_n^2},
\end{split}
\end{equation}
or 
\begin{equation}
\begin{split}
  \sigma_n^2 \approx  C_0^{2/3} \frac{C_n^2 L_0^{2/3}}{2} \implies \frac{C_n^2}{2\sigma_n^2} 
  =  \frac{2 C_n^2}{2 C_0^{2/3} C_n^{2} L_0^{2/3}} = \frac{1}{ C_0^{2/3} L_0^{2/3}} 
\end{split}
\end{equation}
and
\begin{equation}
  C_{cor}(\|\vect{\dim{r}}_1-\vect{\dim{r}}_{2}\|) = 1 - \frac{1}{C_0^{2/3}}
    \begin{cases}
        \frac{1}{(L_0l_{0}^2)^{2/3}}\|\vect{\dim{r}}_1-\vect{\dim{r}}_{2}\|^{2} & 0 < \|\vect{\dim{r}}_1-\vect{\dim{r}}_{2}\| \leq l_{0} \\
        \frac{1}{L_0^{2/3}} \|\vect{\dim{r}}_1-\vect{\dim{r}}_{2}\|^{2/3} &  l_{0}  \leq \|\vect{\dim{r}}_1-\vect{\dim{r}}_{2}\| \leq  L_{0}  \\
   \end{cases}
\end{equation}

\subsubsection{Atmospheric Scales}
\begin{tabular}{ |p{5cm}|p{3cm}|p{4.5cm}|p{3cm}| }
\hline
\textbf{Quantity} & \textbf{Symbol} & \textbf{Characteristic Value(s)} & \textbf{Ref.} \\
\hline
\hline
Wavelength               & $\lambda $    & $10^{-6} $ (meters)                    
& \cite{} \\
\hline
Wavelength               & $k= \frac{2\pi}{\lambda}$      & $10^{6} $ (meters$^{-1}$)    & \\               
\hline
Kolmogorov Inner-scale   & $l_0$         & $10^{-3}-10^{-2}$ (meters)             &  Wheelon, 2005 ,page 32 \\
\hline 
Kolmogorov Outer-scale   & $L_0$         & $10-10^{3}$ (meters)                   &  Wheelon, 2005 ,page 32 \\
\hline
Index Structure Constant & $C_n^2$       & $10^{-20}-10^{-12}$ (meters$^{-2/3}$)  &  Wheelon, 2005 ,page 32\\
\hline
Index Variance           & $\sigma_n^2$  & $C_0^{2/3} (10^{-18}-10^{-10})$        & --\\
\hline
Aperture Diameter        & $D_{ap}$      & $10^{-2}-10^{0}$ (meters)              & \\
\hline
Propagation Distance     & $L$           & $10^{3}-10^{4}$ (meters)               & \\
\hline
\end{tabular}









